\section{Functoriality in the field: what extends along embeddings, and what
does not}
\label{sec:functoriality}

\paragraph{Recollection.} A phase datum for $\kappa=\mathbb{F}_{p^m}$ is a
nontrivial additive character $\psi:(\kappa,+)\to\mathbb{C}^\times$, and the
trace-normalized family is
$\psi_\zeta(z) := \zeta^{\mathrm{Tr}_{\kappa/\mathbb{F}_p}(z)}$ for $\zeta$ a
primitive $p$-th root of unity and
$\mathrm{Tr}_{\kappa/\mathbb{F}_p}(z):=z+z^p+\cdots+z^{p^{m-1}}$
(Definition~\ref{def:phase-datum}); by Theorem~\ref{thm:wh-choice}(a),
$X(\kappa)$ is a $\kappa^\times$-torsor, so every $\psi\in X(\kappa)$ is
$\psi_c := \psi_\zeta(c\,\cdot)$ for a unique $c\in\kappa^\times$.
\provenance{D3}{---}{---}{---}

Earlier sections fixed a single finite field $\kappa$ throughout. This
section asks how the construction behaves as $\kappa$ varies along field
embeddings: whether the trace-normalized family is itself compatible with
restriction, whether the assignment $\kappa\mapsto A_{\psi,\beta_0}(V(\kappa))$
is a functor, and — the question this campaign's earlier working hypothesis
answered wrongly — whether a single compatible choice of character can be
made once and for all, for every finite field simultaneously. It cannot; the
precise obstruction is Frobenius, and the precise set of compatible choices
that \emph{does} exist is named exactly.

\statusProved
\begin{theorem}[The trace-normalized family, restricted along an extension]
\label{thm:wh-funct-a}
Let $\kappa=\mathbb{F}_{p^m}\subseteq\kappa'=\mathbb{F}_{p^{mn}}$. Then
$\psi'_\zeta|_\kappa = \psi_{\zeta^n}$ (both sides characters of $\kappa$,
computed with the same $\zeta$): the restriction of $\kappa'$'s
trace-normalized character to $\kappa$ is trivial exactly when $p\mid n$
(smallest instance: $\mathbb{F}_2\subset\mathbb{F}_4$).
\end{theorem}

\begin{scope}
This theorem is about the trace-normalized family specifically — a
particular, and particularly natural, choice inside $X(\kappa)$ for every
$\kappa$ — and asserts nothing about arbitrary families of characters; the
statement that quantifies over \emph{all} compatible systems is
Theorem~\ref{thm:wh-funct-c} below.
\end{scope}

\provenance{D3, WH-FUNCT-a}{theory/wh-kappa-choice.md \S10 \textit{step} $\langle 1\rangle 1$}{theory/checks/funct\_sections.py (F1)}{theory/verdicts/wh-kappa-adjudication.md}

\begin{proof}
For $x\in\kappa$, $x^{p^m}=x$ (elements of $\mathbb{F}_{p^m}$ are exactly
the fixed points of the $m$-th power of Frobenius), so the $mn$ summands of
$\mathrm{Tr}_{\kappa'/\mathbb{F}_p}(x)=\sum_{k=0}^{mn-1}x^{p^k}$ repeat with
period $m$: grouping into $n$ blocks of $m$ consecutive terms, each block
sums to $\mathrm{Tr}_{\kappa/\mathbb{F}_p}(x)$ (using
$x^{p^{k+m}}=(x^{p^m})^{p^k}=x^{p^k}$), so
$\mathrm{Tr}_{\kappa'/\mathbb{F}_p}|_\kappa = n\cdot\mathrm{Tr}_{\kappa/\mathbb{F}_p}$.
Hence $\psi'_\zeta(x) = \zeta^{\mathrm{Tr}_{\kappa'/\mathbb{F}_p}(x)} = \zeta^{n\cdot\mathrm{Tr}_{\kappa/\mathbb{F}_p}(x)}
= (\zeta^n)^{\mathrm{Tr}_{\kappa/\mathbb{F}_p}(x)} = \psi_{\zeta^n}(x)$, and
$\zeta^n=1$ (so $\psi_{\zeta^n}\equiv 1$) exactly when $p\mid n$, since
$\zeta$ has exact order $p$.
\end{proof}

\statusProved
\begin{theorem}[The reference algebra is functorial in character-compatible
embeddings]
\label{thm:wh-funct-b}
Let $\mathcal{FF}^\pm$ be the category with objects the pairs $(\kappa,\psi)$,
$\psi\in X(\kappa)$, and morphisms $(\kappa,\psi)\to(\kappa',\psi')$ the
field embeddings $\iota:\kappa\hookrightarrow\kappa'$ with
$\psi'\circ\iota=\psi$ (\emph{character-compatible} embeddings). Then
\[
(\kappa,\psi) \ \longmapsto\ A_{\psi,\beta_0}(V(\kappa))
\]
is a functor from $\mathcal{FF}^\pm$ to the category of unital
$\mathbb{C}$-algebras and injective algebra homomorphisms, respecting
composition and identities.
\end{theorem}

\begin{scope}
This functor uses the reference cocycle $\beta_0$ throughout, on every
object; it is not claimed that the resulting assignment is independent of
this choice, nor that an analogous functor exists carrying a varying
$\beta_\kappa$ compatibly along embeddings — the reference cocycle is used
here because it is a fixed, explicit stipulation
(Definition~\ref{def:admissible-cocycles-restated}), not because it is
canonical (Theorem~\ref{thm:wh-beta-c} of \S\ref{sec:polarizing-cocycle}
shows at $p=2$ it is not). Nor is this a statement about the quantum
\emph{system} — the Hilbert space or the Heisenberg group — functorially in
$\kappa$, only about the algebra.
\end{scope}

\provenance{D2, D4, WH-ALG, WH-FUNCT-b}{theory/wh-kappa-choice.md \S10 \textit{step} $\langle 1\rangle 2$}{theory/checks/funct\_sections.py (F5)}{theory/verdicts/wh-kappa-adjudication.md}

\begin{proof}
For a character-compatible embedding $\iota:(\kappa,\psi)\to(\kappa',\psi')$,
put $V(\iota)(a,b):=(\iota a,\iota b) : V(\kappa)\to V(\kappa')$, additive
since $\iota$ is a field homomorphism. Since $\iota$ is multiplicative,
$\beta_0\big(V(\iota)v,V(\iota)v'\big) = \iota(a)\,\iota(b') = \iota(ab') = \iota\big(\beta_0(v,v')\big)$.
Define $W_\beta(v)\mapsto W'_{\beta_0}(V(\iota)v)$ on basis elements of
$A_{\psi,\beta_0}(V(\kappa))$. This sends distinct basis elements to
distinct basis elements ($V(\iota)$ is injective, $\iota$ being a field
embedding), and is multiplicative: the image of
$W_{\beta_0}(v)W_{\beta_0}(v')=\psi(\beta_0(v,v'))W_{\beta_0}(v+v')$ is
$\psi(\beta_0(v,v'))\,W'_{\beta_0}(V(\iota)(v+v'))$, while the product of
the images is
$\psi'\big(\beta_0(V(\iota)v,V(\iota)v')\big)\,W'_{\beta_0}(V(\iota)(v+v'))
= \psi'\big(\iota(\beta_0(v,v'))\big)\,W'_{\beta_0}(V(\iota)(v+v'))
= \psi(\beta_0(v,v'))\,W'_{\beta_0}(V(\iota)(v+v'))$,
using $\psi'\circ\iota=\psi$ (the compatibility condition). So the map is a
unital, nonzero algebra homomorphism, hence injective by simplicity of the
source (Theorem~\ref{thm:wh-alg}). It respects composition and identities
because $\iota\mapsto V(\iota)$ does: $V(\iota'\circ\iota)=V(\iota')\circ V(\iota)$
and $V(\mathrm{id})=\mathrm{id}$, directly from the definition of $V(\iota)$.
\end{proof}

\subsection{No compatible choice of character over every finite field}

\statusRefuted
\begin{theorem}[Refuted: a Frobenius-compatible character exists for every
finite field]
\label{thm:wh-funct-b-sec}
\emph{(Claimed, and refuted.)} A section of the forgetful functor
$\mathcal{FF}^\pm \to \{\text{finite fields}\}$ exists: a choice
$\psi_\kappa\in X(\kappa)$ for every finite field $\kappa$ such that
$\iota:\kappa\to\kappa'$ character-compatible (i.e.\
$\psi_{\kappa'}\circ\iota=\psi_\kappa$) for \emph{every} field embedding
$\iota$ — in particular, for the Frobenius self-embedding of every $\kappa$.
\end{theorem}

\begin{scope}
This row is kept, refuted, rather than deleted, per this campaign's status
discipline: a refutation carrying the surviving weaker statement is a
result, not a gap to be erased. The surviving statement — exactly which
obstruction blocks a section, and exactly which weaker compatible systems
\emph{do} exist — is Theorem~\ref{thm:wh-funct-c} immediately below and
Theorem~\ref{thm:wh-funct-d} thereafter. Nothing may depend on this row.
\end{scope}

\provenance{WH-FUNCT-b-SEC}{---}{theory/checks/funct\_sections.py (F2,F3)}{theory/verdicts/wh-kappa-adjudication.md}

\begin{quote}
\emph{Refutation.} Frobenius $x\mapsto x^p$ is a field embedding
$\kappa\to\kappa$, so a section would in particular have to satisfy
$\psi_\kappa\circ\mathrm{Frob}=\psi_\kappa$ for every $\kappa$. Theorem~\ref{thm:wh-funct-c}
below shows this forces $\psi_\kappa=\psi_c$ for some $c$ in the \emph{prime}
field $\mathbb{F}_p$, and that for $\kappa=\mathbb{F}_{p^p}$ every character
of this Frobenius-invariant form restricts \emph{trivially} to
$\mathbb{F}_p$ — contradicting $\psi_{\mathbb{F}_p}\in X(\mathbb{F}_p)$,
which must be nontrivial by definition. So no section exists; the smallest
instance of the obstruction is $\mathbb{F}_2\subset\mathbb{F}_4$. A
previously considered route to a section, by induction along the tower
$\mathbb{F}_{p^{n!}}$, is insufficient in any case: a character of
$\mathbb{F}_{p^{n!}}$ compatible at that one stage can still restrict
trivially to an intermediate field, so no such induction can repair the
obstruction just exhibited.
\end{quote}

\statusProved
\begin{theorem}[The Frobenius obstruction, stated exactly]
\label{thm:wh-funct-c}
There is no choice of $\psi_\kappa\in X(\kappa)$, one for each finite field
$\kappa$, compatible with \emph{all} field embeddings. Precisely: writing
$\psi_c := \psi_\zeta(c\,\cdot)$ for $c\in\kappa$ (Theorem~\ref{thm:wh-choice}(a)),
a character $\psi_c$ is invariant under the Frobenius self-embedding of
$\kappa=\mathbb{F}_{p^m}$ if and only if $c\in\mathbb{F}_p$; and for
$\kappa=\mathbb{F}_{p^p}$, every $\mathbb{F}_p$-valued $c$ gives a character
$\psi_c$ that restricts \emph{trivially} to the prime field
$\mathbb{F}_p\subset\mathbb{F}_{p^p}$. The smallest instance is
$\mathbb{F}_2\subset\mathbb{F}_4$.
\end{theorem}

\begin{scope}
This theorem exhibits the precise obstruction at the specific extension
degree $m=p$; it does not enumerate every extension degree at which some
weaker compatibility fails, and it does not address (this is
Theorem~\ref{thm:wh-funct-d}) which \emph{restricted} systems of compatible
characters — over a sub-poset of embeddings, or with the ``for every
embedding'' requirement relaxed — do exist.
\end{scope}

\provenance{D3, WH-FUNCT-a, WH-FUNCT-c}{theory/wh-kappa-choice.md \S10 \textit{step} $\langle 1\rangle 3$}{theory/checks/funct\_sections.py (F2,F3)}{theory/verdicts/wh-kappa-adjudication.md}

\begin{proof}
\emph{Frobenius-invariance forces $c\in\mathbb{F}_p$.} Write
$d:=c^{1/p}$ (the unique $p$-th root, Frobenius being bijective on the
finite field $\kappa$). Since
$\mathrm{Tr}_{\kappa/\mathbb{F}_p}(y^p)=\mathrm{Tr}_{\kappa/\mathbb{F}_p}(y)$
for every $y$ (Frobenius cyclically permutes the $m$ summands of the trace,
using $y^{p^m}=y$), we get
$\psi_c(x^p) = \zeta^{\mathrm{Tr}(cx^p)} = \zeta^{\mathrm{Tr}((dx)^p)} = \zeta^{\mathrm{Tr}(dx)} = \psi_d(x)$,
i.e.\ $\psi_c\circ\mathrm{Frob}=\psi_d$. Since $c\mapsto\psi_c$ is injective
on $\kappa$ (if $\psi_c=\psi_{c'}$ then $\mathrm{Tr}((c-c')x)\equiv 0$ for
all $x$, forcing $c=c'$ by the same nondegeneracy of the trace pairing used
in Theorem~\ref{thm:wh-choice}(a)'s proof), $\psi_c$ is Frobenius-invariant
iff $d=c$; since $d^p=c$ by definition of $d$, substituting $d=c$ gives
$c^p=c$, and conversely $c^p=c$ gives $d=c^{1/p}=c$ (Frobenius being
bijective). So the condition is exactly $c^p=c$, i.e.\ $c$ is a root of
$X^p-X$, i.e.\
(this polynomial having at most $p$ roots in a field, all $p$ elements of
$\mathbb{F}_p$ among them by Fermat's little theorem for finite fields)
$c\in\mathbb{F}_p$.

\emph{Trivial restriction at $\kappa=\mathbb{F}_{p^p}$.} Here $m=p$. For
$c\in\mathbb{F}_p$ and $x\in\mathbb{F}_p\subset\kappa$, every summand of
$\mathrm{Tr}_{\kappa/\mathbb{F}_p}(cx)=\sum_{k=0}^{p-1}(cx)^{p^k}$ equals
$cx$ (as $cx\in\mathbb{F}_p$ is already Frobenius-fixed), so
$\mathrm{Tr}_{\kappa/\mathbb{F}_p}(cx)=p\cdot(cx)=0$ in characteristic $p$.
So $\psi_c(x)=\zeta^0=1$ for every $x\in\mathbb{F}_p$: \emph{every}
Frobenius-invariant character of $\mathbb{F}_{p^p}$ — the only kind a
section could assign there — restricts trivially to the prime field,
contradicting the requirement $\psi_{\mathbb{F}_p}\in X(\mathbb{F}_p)$
(nontrivial by Definition~\ref{def:phase-datum}) forced by compatibility
along the inclusion $\mathbb{F}_p\subset\mathbb{F}_{p^p}$. At $p=2$,
$\mathbb{F}_{p^p}=\mathbb{F}_4$: the smallest instance.
\end{proof}

\subsection{The compatible systems that do exist}

\statusProved
\begin{theorem}[Compatible character systems over an algebraic closure]
\label{thm:wh-funct-d}
Fix an algebraic closure $\overline{\mathbb{F}}_p$, and consider the
inclusion poset of its finite subfields. The compatible systems
$(\psi_\kappa)_\kappa$ — one $\psi_\kappa\in X(\kappa)$ for every finite
subfield $\kappa\subset\overline{\mathbb{F}}_p$, with
$\psi_{\kappa'}|_\kappa=\psi_\kappa$ for every inclusion $\kappa\subseteq\kappa'$
of finite subfields — are \textbf{exactly} the restrictions
$\psi_\kappa:=\Psi|_\kappa$ of the additive characters $\Psi$ of
$(\overline{\mathbb{F}}_p,+)$ with $\Psi|_{\mathbb{F}_p}\neq 1$.
\end{theorem}

\begin{scope}
This theorem restricts attention to the inclusion poset of finite subfields
of a \emph{fixed} algebraic closure, respecting compatibility only along
inclusions of subfields, not along all abstract field embeddings between
isomorphic-but-not-identical finite fields; Theorem~\ref{thm:wh-funct-b-sec}
shows the latter, stronger requirement admits no solution at all. It says
nothing about compatible choices of a polarizing cocycle $\beta$ along the
same poset; whether the Arf type of \S\ref{sec:polarizing-cocycle} can be
chosen compatibly there is not resolved by this labbook.
\end{scope}

\provenance{D3, WH-FUNCT-d}{theory/wh-kappa-choice.md \S10 \textit{step} $\langle 1\rangle 4$}{theory/checks/funct\_sections.py (F4)}{theory/verdicts/wh-kappa-adjudication.md}

\begin{proof}
\emph{Restrictions are compatible systems.} Given such a $\Psi$, put
$\psi_\kappa:=\Psi|_\kappa$; compatibility along $\kappa\subseteq\kappa'$
holds by construction ($(\Psi|_{\kappa'})|_\kappa=\Psi|_\kappa$), and
$\psi_\kappa$ is nontrivial for \emph{every} finite subfield $\kappa$: every
such $\kappa$ contains $\mathbb{F}_p$, and $\Psi|_\kappa\equiv 1$ would
force $\Psi|_{\mathbb{F}_p}\equiv 1$ too, contradicting the hypothesis on
$\Psi$.

\emph{Compatible systems arise this way.} Given a compatible system
$(\psi_\kappa)$, define $\Psi(x):=\psi_\kappa(x)$ for any finite subfield
$\kappa\ni x$ (every $x\in\overline{\mathbb{F}}_p$ lies in some finite
subfield, $\overline{\mathbb{F}}_p$ being their union). This is well
defined: if $x$ lies in two finite subfields $\kappa,\kappa'$, both lie
inside a common finite subfield $\kappa''$ (e.g.\ their compositum, again
finite), and compatibility along $\kappa\subseteq\kappa''$ and
$\kappa'\subseteq\kappa''$ gives $\psi_\kappa(x)=\psi_{\kappa''}(x)=\psi_{\kappa'}(x)$.
$\Psi$ is additive: any two elements lie in a common finite subfield, where
$\Psi$ agrees with the additive character $\psi_\kappa$ of that subfield.
And $\Psi|_{\mathbb{F}_p}=\psi_{\mathbb{F}_p}\neq 1$, by
Definition~\ref{def:phase-datum}'s requirement on every $\psi_\kappa\in X(\kappa)$
applied to $\kappa=\mathbb{F}_p$.

\emph{The two constructions are mutually inverse}, and the condition
``nontrivial on every finite subfield'' (needed termwise, since each
$\psi_\kappa\in X(\kappa)$ by hypothesis) is equivalent to ``nontrivial on
the prime field alone'': the forward implication is immediate ($\mathbb{F}_p$
is itself one of the subfields), and the reverse is exactly what the first
paragraph showed.
\end{proof}

\begin{remark}[What this section does not claim]
Theorem~\ref{thm:wh-funct-b} is functoriality of the \emph{algebra} in
character-compatible embeddings, for the fixed reference cocycle
$\beta_0$ — not functoriality of the quantum system (Heisenberg group,
Hilbert space) itself, and not a claim that any compatible choice of
polarizing cocycle along embeddings exists. Whether the Arf type of
Theorem~\ref{thm:wh-beta-type} can be chosen compatibly along field
embeddings at $p=2$ is recorded as an open question by this campaign's
source shard, not resolved here.
\end{remark}
